% ---------------------------------------------------------------------
\chapter{Wie Git denkt: das mentale Modell}
\label{kap:modell}

\begin{lernziele}
\begin{itemize}[nosep]
  \item die vier Orte kennen, zwischen denen Änderungen wandern,
  \item verstehen, was ein Commit technisch ist,
  \item wissen, was Branches, HEAD und Remotes sind,
  \item jeden Git-Befehl einem Weg zwischen den vier Orten zuordnen können.
\end{itemize}
\end{lernziele}

Die meisten Schwierigkeiten mit Git entstehen nicht durch vergessene Befehle, sondern durch
ein unklares Bild davon, \emph{wo} sich eine Änderung gerade befindet. Dieses Kapitel
vermittelt dieses Bild. Es ist das wichtigste Kapitel des Handbuchs.

\section{Die vier Orte}

Eine Änderung durchläuft auf ihrem Weg zu GitHub vier Orte. Jeder Git-Befehl bewegt
Änderungen von einem Ort zum nächsten oder zeigt an, was wo liegt.

\begin{center}
\begin{tikzpicture}[
    ort/.style={draw=hsluBlau, fill=hsluBlau!6, rounded corners=3pt, thick,
                minimum width=3.0cm, minimum height=1.5cm, align=center, font=\small},
    weg/.style={-{Stealth[length=2.5mm]}, thick},
    lab/.style={font=\ttfamily\footnotesize, fill=white, inner sep=1pt}
  ]
  \node[ort] (arbeit) at (0,0)    {\textbf{Arbeitsordner}\\deine Dateien};
  \node[ort] (vor)    at (3.9,0)  {\textbf{Staging Area}\\Vormerkung};
  \node[ort] (repo)   at (7.8,0)  {\textbf{lokales Repo}\\Geschichte auf\\deinem PC};
  \node[ort] (gh)     at (11.7,0) {\textbf{Remote}\\GitHub};

  \draw[weg] ([yshift=4mm]arbeit.east) -- node[lab, above=1mm] {add}     ([yshift=4mm]vor.west);
  \draw[weg] ([yshift=-4mm]vor.west)   -- node[lab, below=1mm] {restore} ([yshift=-4mm]arbeit.east);
  \draw[weg] (vor.east)                -- node[lab, above=1mm] {commit}  (repo.west);
  \draw[weg] ([yshift=4mm]repo.east)   -- node[lab, above=1mm] {push}    ([yshift=4mm]gh.west);
  \draw[weg] ([yshift=-4mm]gh.west)    -- node[lab, below=1mm] {pull}    ([yshift=-4mm]repo.east);
\end{tikzpicture}
\end{center}

\begin{description}[style=nextline, leftmargin=1.2cm]
  \item[1. Arbeitsordner \normalfont(\emph{working tree})]
    Die Dateien, die im Explorer und im Editor sichtbar sind. Hier wird gearbeitet. Git
    beobachtet diesen Ordner, verändert aber nichts von sich aus.
  \item[2. Staging Area \normalfont(Vormerkung, auch \emph{index})]
    Eine Zwischenablage für den nächsten Zwischenstand. Man legt dort gezielt die
    Änderungen ab, die zusammengehören.
  \item[3. Lokales Repository \normalfont(\emph{repository}, kurz \emph{Repo})]
    Die gespeicherte Geschichte des Projekts. Sie liegt im versteckten Ordner \cmd{.git}
    innerhalb des Projektordners. Wer diesen Ordner löscht, löscht die Geschichte.
  \item[4. Remote \normalfont(entfernte Kopie)]
    Eine Kopie des Repositorys auf einem anderen Rechner, meist auf GitHub. Sie dient als
    Sicherung und als gemeinsamer Treffpunkt für alle Beteiligten.
\end{description}

\subsection*{Warum gibt es die Staging Area?}

Die Staging Area irritiert Einsteiger am meisten: Warum nicht direkt speichern? Der Grund
ist, dass bei der Arbeit oft mehrere Dinge gleichzeitig entstehen. Angenommen, an einem
Abend wurde eine Übungsaufgabe gelöst \emph{und} nebenbei ein Tippfehler in der
Dokumentation korrigiert. Das sind zwei verschiedene Anliegen. Mit der Staging Area lassen
sie sich als zwei getrennte, sauber beschriebene Zwischenstände festhalten.

Ein Vergleich: Die Staging Area ist wie ein Einkaufskorb. Man legt hinein, was zusammengehört,
und erst an der Kasse, beim \emph{Commit}, wird verbucht.

\section{Der Commit: ein Foto des Projekts}

\begin{definition}{Commit}
Ein \begriff{Commit} ist ein gespeicherter Zwischenstand des gesamten Projekts. Er enthält
\begin{itemize}[nosep]
  \item den Inhalt aller Dateien zu diesem Zeitpunkt,
  \item Autor und Zeitpunkt,
  \item eine Beschreibung, die \emph{Commit-Nachricht},
  \item einen Verweis auf den vorherigen Commit.
\end{itemize}
\end{definition}

Ein Commit ist also kein „Änderungsprotokoll“, sondern ein vollständiges \textbf{Foto} des
Projekts. Weil jeder Commit auf seinen Vorgänger verweist, entsteht eine Kette, die
\begriff{Geschichte} (\emph{history}):

\begin{center}
\begin{tikzpicture}[
    commit/.style={circle, draw=hsluBlau, fill=hsluBlau!10, thick, minimum size=9mm,
                   font=\ttfamily\scriptsize},
    vor/.style={-{Stealth[length=2mm]}, thick, gray}
  ]
  \node[commit] (c1) at (0,0)   {98bc};
  \node[commit] (c2) at (2.2,0) {19ce};
  \node[commit] (c3) at (4.4,0) {3b1a};
  \node[commit] (c4) at (6.6,0) {7f02};
  \draw[vor] (c2) -- (c1);
  \draw[vor] (c3) -- (c2);
  \draw[vor] (c4) -- (c3);
  \node[font=\small, below=2mm of c1] {ältester};
  \node[font=\small, below=2mm of c4] {neuester};
\end{tikzpicture}
\end{center}

Die Pfeile zeigen rückwärts: Jeder Commit kennt seinen Vorgänger, nicht seinen Nachfolger.

\subsection*{Der Hash: die Adresse eines Commits}

Jeder Commit erhält eine eindeutige Kennung, den \begriff{Hash}, zum Beispiel
\cmd{3b1a03a5c2e0...}. Er besteht aus 40 Zeichen. Im Alltag genügen die ersten sieben, etwa
\cmd{3b1a03a}. Mit dem Hash lässt sich jeder Commit genau ansprechen, zum Beispiel um ihn
anzuzeigen oder rückgängig zu machen.

\begin{hintergrund}
Der Hash wird aus dem Inhalt des Commits berechnet, also aus Dateien, Autor, Zeitpunkt,
Nachricht \emph{und dem Hash des Vorgängers}. Ändert sich auch nur ein Zeichen irgendwo in
der Geschichte, ändern sich alle nachfolgenden Hashes. Darum ist die Geschichte in Git
praktisch fälschungssicher, und darum lassen sich bereits geteilte Commits nicht einfach
„umschreiben“, ohne dass alle anderen es merken (siehe Kapitel~\ref{kap:fehler}).
\end{hintergrund}

\section{Branches: parallele Arbeitslinien}

\begin{definition}{Branch}
Ein \begriff{Branch} (deutsch: Zweig) ist eine eigene Arbeitslinie in der Geschichte.
Technisch ist er nur ein beweglicher Zeiger auf einen Commit. Mit jedem neuen Commit rückt
der Zeiger des aktuellen Branches automatisch nach vorn.
\end{definition}

Jedes Repository hat einen Hauptbranch, heute meist \cmd{main} genannt (früher
\cmd{master}). Er enthält den gültigen, funktionierenden Stand. Neue Arbeit geschieht auf
einem eigenen Branch, der von \cmd{main} abzweigt:

\begin{center}
\begin{tikzpicture}[
    commit/.style={circle, draw=hsluBlau, fill=hsluBlau!10, thick, minimum size=7mm},
    neu/.style={circle, draw=bspFarbe, fill=bspFarbe!15, thick, minimum size=7mm},
    vor/.style={-{Stealth[length=2mm]}, thick, gray},
    zeiger/.style={rounded corners=2pt, font=\ttfamily\small, inner sep=3pt, text=white}
  ]
  \node[commit] (c1) at (0,0) {};
  \node[commit] (c2) at (1.8,0) {};
  \node[commit] (c3) at (3.6,0) {};
  \node[neu]    (d1) at (5.4,1.2) {};
  \node[neu]    (d2) at (7.2,1.2) {};
  \draw[vor] (c2) -- (c1);
  \draw[vor] (c3) -- (c2);
  \draw[vor] (d1) -- (c3);
  \draw[vor] (d2) -- (d1);
  \node[zeiger, fill=hsluBlau, below=4mm of c3] (m) {main};
  \draw[-{Stealth}, hsluBlau, thick] (m) -- (c3);
  \node[zeiger, fill=bspFarbe, above=4mm of d2] (a) {ads};
  \draw[-{Stealth}, bspFarbe, thick] (a) -- (d2);
  \node[font=\small\ttfamily, right=3mm of a, text=gray] (h) {HEAD};
  \draw[-{Stealth}, gray, thick] (h) -- (a);
\end{tikzpicture}
\end{center}

Hier zeigt \cmd{main} weiterhin auf den letzten gültigen Stand. Auf dem Branch \cmd{ads}
sind zwei neue Commits entstanden. \cmd{main} bleibt davon unberührt, bis die Arbeit
geprüft und übernommen wird. Dieses Übernehmen heisst \begriff{Merge} (zusammenführen).

\begin{definition}{HEAD}
\begriff{HEAD} ist der Zeiger auf „hier stehe ich gerade“. Er zeigt auf den aktuellen
Branch. Wechselt man den Branch, wandert HEAD mit, und Git passt die Dateien im
Arbeitsordner an den Stand dieses Branches an.
\end{definition}

\begin{praxis}
Weil ein Branch nur ein Zeiger ist, kostet er praktisch nichts: Anlegen, Wechseln und
Löschen dauert Sekundenbruchteile. Darum arbeiten Teams mit sehr vielen kurzlebigen
Branches, typischerweise einem pro Aufgabe oder Fehlerbehebung. Ein Branch lebt oft nur
wenige Stunden oder Tage, bis er in \cmd{main} übernommen und wieder gelöscht wird.
\end{praxis}

\section{Remotes: die Kopie im Internet}

Ein \begriff{Remote} ist ein anderes Repository, mit dem das eigene synchronisiert wird.
Beim Klonen von GitHub erhält es automatisch den Namen \cmd{origin} (Ursprung). Mit zwei
Befehlen wird abgeglichen:

\begin{itemize}
  \item \cmd{git push} schickt eigene Commits zum Remote.
  \item \cmd{git pull} holt fremde Commits vom Remote.
\end{itemize}

Git merkt sich zusätzlich, auf welchem Stand das Remote beim letzten Abgleich war. Dieser
gespeicherte Stand heisst zum Beispiel \cmd{origin/main}. Daher stammen Meldungen wie
\emph{„Your branch is ahead of \textquotesingle{}origin/main\textquotesingle{} by 1 commit“}:
Der eigene \cmd{main} hat einen Commit mehr als der zuletzt bekannte Stand auf GitHub.

\begin{zusammenfassung}
\begin{itemize}[nosep]
  \item Änderungen wandern: Arbeitsordner $\rightarrow$ Staging Area $\rightarrow$
        lokales Repo $\rightarrow$ Remote.
  \item Ein Commit ist ein vollständiges Foto des Projekts mit Autor, Zeit, Nachricht und
        Verweis auf den Vorgänger. Seine Adresse ist der Hash.
  \item Ein Branch ist ein beweglicher Zeiger auf einen Commit. \cmd{main} ist der gültige Stand.
  \item HEAD zeigt, wo man gerade steht. \cmd{origin} ist die Kopie auf GitHub.
\end{itemize}
\end{zusammenfassung}
